\section{Methodology}\label{sec:methodology}

This section presents the system in two equally weighted halves: detection
and federation.

\subsection{System and Problem Formulation}

The detector is embedded in the three-tier system of Fig.~\ref{fig:f1}: a
\emph{browser tier} that builds the feature vector and infers locally, a
\emph{platform tier} that holds identity and the aggregation core and never
receives holdings, and a \emph{model tier} that is the single authoritative
trainer. Two versioned contracts hold the tiers together --- a feature/scaler
specification and a set of parity fixtures --- detailed in
Section~\ref{sec:methodology}-E, since they are preconditions for aggregation
rather than incidental engineering. Only a twelve-element coefficient vector,
a scalar intercept and a sample count ever cross a tier boundary.

Let $x\in\mathbb{R}^{12}$ describe a portfolio and its market context on a
decision day. Rebalancing-necessity detection is the binary task
\begin{equation}\label{eq:e1}
\hat{y} = g\bigl(f(x;\theta)\bigr),\qquad y\in\{0,1\},
\end{equation}
where $y=1$ denotes that rebalancing is warranted. This models the
\emph{decision}, not the trade schedule implementing it, which is a separate
deterministic engine outside this scope.

\subsection{Features, Dataset and Labels}

The contract fixes twelve quantities in order: seven portfolio descriptors, one
elapsed-time descriptor and four market descriptors, defined together with
their distribution and univariate association with the label in
Table~\ref{tab:t2}. The vector is assembled in that fixed order and
standardized by shared constants,
\begin{equation}\label{eq:e16}
z_i = (x_i - \mu_i)/\sigma_i \ \ \text{if } \sigma_i \geq 10^{-9}, \ \ \text{else } 0,
\end{equation}
part of the versioned contract rather than fitted per run, since federated
averaging is coherent only when every participant occupies the same input
space. This carries a disclosed defect: the shipped constants derive from a
legacy 5{,}000-row dataset rather than the 30{,}000-row training split used
here (largest discrepancy on \texttt{num\_stocks}, $|\Delta\mu|{=}3.9504$); the
defect is disclosed rather than repaired, and
Section~\ref{sec:implementation} quantifies its cost at $+0.0001$ test AUC.

No public dataset pairs portfolio states with realized rebalancing outcomes,
so the corpus is constructed --- generated data is established practice in
financial machine learning where labelled outcomes are unobtainable
\cite{mashrur2020}. The generator emits $30{,}000$ states under seed
$20260804$; each row is one simulated decision on a real trading day. The
\emph{market leg is real}: NIFTY~50 daily closes over $3{,}706$ sessions
(2011-08-01 to 2026-07-31) supply every return and market feature, while only
the cross-section of holdings is generated from one of six basket profiles
with log-normal target weights and a persistent per-symbol factor loading
$\beta_i\in[0.55,1.55]$. The book drifts over $5$--$170$ trading days by a
real market return plus a systematic, sector and idiosyncratic shock, and
forward returns combine the real forward index move with a regime-dependent
reversion term, stronger under volatility/drawdown stress and weaker in calm
uptrends~\cite{jegadeesh1993}. Persistent loadings keep the market factor from
averaging out, regime-dependence keeps returns from collapsing into a function
of weights alone, and drift is \emph{caused} by elapsed time, so elapsed time
reaches the label only along that causal path; exact constants are in the
repository alongside the reproduction script.

The label is the outcome of a forward, net-of-cost backtest rather than a
rule. Let $t$ be target weights, $w$ the drifted current weights and $f$ the
realized forward returns. The investor compares rebalancing against holding
under a mean-variance objective with transaction costs
\cite{markowitz1952,sharpe1966}, and the label is the sign of the resulting net
certainty-equivalent gain,
\begin{equation}\label{eq:e23}
\begin{aligned}
\Delta\mathrm{CE}(w,t,f) ={}& (t-w)^{\!\top} f
+ A\bigl[\mathrm{Var}(w)-\mathrm{Var}(t)\bigr] \\
&- c\,\lVert t-w \rVert_1,\qquad
y = \mathbf{1}\{\Delta\mathrm{CE} > 0\},
\end{aligned}
\end{equation}
with risk aversion $A{=}45.0$ and round-trip cost $c{=}0.0035$ ($35$ basis
points): a \emph{return gain}, a \emph{risk gain} priced through ex-ante
variance, and an \emph{implementation cost}. Because cost scales with
turnover, \eqref{eq:e23} reproduces no-trade-band behaviour without any band
being coded. Over $30{,}000$ rows it yields a positive rate of $0.5009$
obtained structurally, without resampling. The elapsed-time feature $\tau$ is
deliberately absent from \eqref{eq:e23}; it influences the label only through
the drift it causes, removing the time-collinearity defect of the predecessor
heuristic, which instead scored a sigmoid combination of a logistic function
of $\tau$ and a tiered adjustment on drift, concentration, drawdown, trend and
volatility directly from the feature vector. Because that score is a
deterministic function of the same features the classifier receives, a model
trained on its output recovers the rule rather than the decision; it is used
here solely as a negative control, never as the reported detector's label
source.

\subsection{Detection Model and Training}

The detector is a binary logistic regression on the standardized vector of
\eqref{eq:e16},
\begin{equation}\label{eq:e25}
P(y=1\mid x) = \sigma\bigl(\langle \mathbf{w}, z \rangle + b\bigr), \quad
\hat{y} = \mathbf{1}\{P \geq 0.5\},
\end{equation}
fitted by $L_2$-regularized maximum likelihood ($C{=}1$, lbfgs,
\texttt{max\_iter}$=1000$). The $30{,}000$ rows split
$60/20/20$ by source row order into $18{,}000$/$6{,}000$/$6{,}000$, described as
row-order rather than temporal since the corpus carries no market-time index.

The linear model is a deployment choice, not an accuracy choice: its
coefficients are readable log-odds contributions (Table~\ref{tab:t2}) and are
exactly the object federated averaging operates on --- a weight vector plus
intercept with no architectural ambiguity, unlike a tree ensemble, which
admits no meaningful coordinate-wise mean. Identical arithmetic in both
runtimes is what makes the parity contract testable to $10^{-12}$, and the
objective yields calibrated probabilities for an interface that presents a
probability of need rather than a binary instruction. The accuracy price is
quantified in Section~\ref{sec:implementation}.

\subsection{Federated Protocol and Aggregation}

The parameters of \eqref{eq:e25} are improved across accounts without any
account disclosing its holdings. Fig.~\ref{fig:f12} shows the round lifecycle:
a model is seeded from the reference corpus, distributed to authenticated
clients that refuse a contract mismatch, fitted locally in the browser
(warm-started from the global parameters under Adam, lr $0.05$, batch $32$,
so the local optimum is a displacement of the global model), and uploaded ---
twelve coefficients, an intercept, the sample count $n_k$ and the contract
versions, nothing else --- validated for shape, finiteness and contract match
before storage (mismatches rejected as HTTP 409). This scope is stated
without ambiguity: this path is \textbf{implemented, validated and
unit-tested, but disabled by default in the shipped product and not
experimentally evaluated in this paper} --- no federated accuracy, convergence
or robustness result is claimed, and every number in
Section~\ref{sec:implementation} comes from centralized training. Transmitting
updates instead of holdings is data minimisation, not a privacy guarantee;
neither differential privacy nor secure aggregation is implemented. Both
statements are repeated once in Section~\ref{sec:conclusion}.

Aggregation operates on displacements. Let $W_{\mathrm{g}}$ be the active
global model; each admitted update is expressed as a delta and clipped in
$L_2$ norm,
\begin{equation}\label{eq:e28}
\Delta_k = W_k - W_{\mathrm{g}}, \qquad
\tilde{W}_k = W_{\mathrm{g}} + \Delta_k \min\!\left(1,\ \tfrac{L}{\lVert \Delta_k \rVert_2}\right),
\end{equation}
with $L{=}5.0$, then combined by a sample-weighted mean~\cite{mcmahan2017},
\begin{equation}\label{eq:e29}
\bar{W} = \sum_k \frac{n_k}{\sum_j n_j}\, \tilde{W}_k .
\end{equation}
Weighting is not cosmetic: two participants contributing $0.4$ on $120$
samples and $0.8$ on $360$ samples aggregate to $(120\cdot0.4+360\cdot0.8)/480=0.7$
under \eqref{eq:e29}, against $0.6$ unweighted. A round executes only under
\begin{equation}\label{eq:e30}
\bigl|\{k : k \text{ valid}\}\bigr| \geq 2 \ \wedge \ \text{contract versions current},
\end{equation}
since a single participant would publish its local model as the global one,
and mixed contracts would average coefficients defined over different feature
or scaler spaces into a vector that is arithmetically well formed and
semantically meaningless. Heterogeneous client distributions are a known
source of degradation for this family of algorithms
\cite{li2020challenges,li2020noniid}; which regime this system falls into is an
open question, not a claim made here.

\subsection{Robustness and Parity Contract}

The assumed adversary is a participating account, or colluding accounts,
attempting to move the global model; the design assumes neither honest
participants nor cryptographic identity, so it bounds what an identity is
worth rather than preventing Sybil creation: clipping in~\eqref{eq:e28} bounds
any single displacement to $L{=}5.0$ regardless of submitted magnitude, a cap
of $n_k\leq10{,}000$ bounds the weight one identity can claim
in~\eqref{eq:e29} (since $n_k$ is self-reported and unverifiable without
seeing data the design is built not to see), and an hourly throttle bounds
both submission rate and reconnaissance. Aggregation is
administrator-authenticated and lock-protected, activation is
push-before-persist and atomic with the ledger, and the append-only ledger
makes recovery a forward operation; residual weaknesses (a single static
secret, cleartext transport, non-cryptographic identity) are recorded rather
than mitigated, and discussed in Section~\ref{sec:conclusion}.

Federated averaging across two runtimes presupposes they compute the same
function. Two versioned artefacts enforce this: \texttt{feature-spec.json}
fixes the feature order and constants of~\eqref{eq:e16}; \texttt{parity-fixtures.json}
pins five reference vectors and their expected standardized forms, logits and
probabilities, asserted by both test suites at $10^{-12}$ (float64) and
$10^{-5}$ (browser float32). The platform refuses to store or aggregate any
row whose contract triple differs from the active one (HTTP 409) rather than
risk an aggregate silently polluted by incompatible participants --- which is
also what lets the fixed-constant standardization of~\eqref{eq:e16} stand
despite its disclosed scaler-provenance defect, since a scaler refitted
per-client would invalidate~\eqref{eq:e29} outright.
